\chapter{Conclusion}
\label{ch:conclusion}

This thesis set out to build an autonomous trading agent on a standardised tool protocol such
that look-ahead bias is impossible by construction, and to evaluate the quality of its
reasoning --- not merely its returns --- rigorously and by market regime.

The methodological contribution is the \emph{\mcp-as-time-machine} design. A single simulation
clock exposes \tnow; every Model Context Protocol data tool refuses, at read time, to return
any datum dated strictly after it; and the agent calls the identical tools in backtest and in
live operation, so the two are exactly equivalent and look-ahead is structurally impossible
rather than merely discouraged (\cref{inv:nolookahead}). This guarantee is pinned by tests
written before the data layer and falsified by mutation testing, and it is complemented by
explicit treatment of the leakage channels a temporal filter does not close --- adjusted-price
leakage and weight-memory contamination --- the latter documented honestly as a residual
threat rather than overclaimed.

The evaluation contribution is a regime-segmented measurement of reasoning quality:
faithfulness (the knowledge--action gap), measured by a blind, policy-aware, temperature-zero
LLM judge and anchored to blind human annotation; grounding, computed deterministically against
recorded tool outputs; and a sophistication rubric modelled on KellyBench. The combination of
near-perfect grounding with a regime-dependent faithfulness gap is the central qualitative
finding, and the diagnosis of that gap as an exit-discipline problem --- visible only once the
judge was made policy-aware and pre-registered --- is itself a methodological result.

The single-agent system is the primary, canonical result; the multi-agent Portfolio Manager is
reported honestly as a bonus that tests whether explicit deliberation changes behaviour. This
hierarchy is deliberate and reflects the thesis's governing principle: the contribution is the
\emph{harness that measures the knowledge--action gap cleanly for tool-using agents}, not the
number of agents in it. The rigour narrative --- a chain of real bugs caught, diagnosed, fixed,
and closed by regression tests --- is offered as the standing evidence that the instruments do
what they claim.

The single-agent claims are confirmed on the full two-year canonical run: grounding is
near-perfect across regimes, faithfulness rises substantially under the pre-registered
policy-aware judge yet remains lowest in the bull regime, and the dominant residual gap is a
specific exit-discipline failure validated against human annotation. The later one-year
common-reference run resolves the two drafting incoherences: single agent, PM, and baselines
share the same window, tickers, warm-up, and $10$\,bps transaction costs. On that clean
benchmark the PM extension is leak-free and evidence-grounded, but its Sharpe confidence
interval crosses zero and its error rate is non-zero, so it remains an architecture extension
rather than a performance claim. The design, the protocol, and their justification do not
depend on asserting PM outperformance.
